\documentclass[11pt,letterpaper,english]{article}
\usepackage{microtype}
\usepackage{titling}
\usepackage{amsmath} 
\usepackage{amsthm}
\usepackage{mathrsfs}
\usepackage{caption}
\usepackage{bibunits}
\usepackage{fontawesome}
\usepackage[font=small,labelfont=bf]{caption}
\usepackage{subcaption}
\usepackage{setspace}
\usepackage{amssymb}
\usepackage{multirow}
\usepackage{rotating}
\usepackage{array}
\usepackage{multibib}
\usepackage{mathtools}
\usepackage{siunitx}
\usepackage{colortbl}
\usepackage{float}
\usepackage[flushleft]{threeparttable}
\usepackage{algorithm}
\usepackage{multirow, makecell}
\makeatletter
\renewcommand{\ALG@name}{Pseudo-code}
\makeatother
\usepackage{algorithmic}
\usepackage{xr}
\usepackage{color}
\DeclareMathAlphabet{\mathcalligra}{T1}{calligra}{m}{n}
\newcommand{\Var}{\mathrm{Var}}
\DeclareMathOperator{\E}{\mathbb{E}}
\DeclareMathOperator{\Tr}{Tr}
\usepackage[sectionbib,round]{natbib}
\setlength{\bibsep}{0.0pt}
\setlength{\footskip}{30pt}
\newtheorem{propn}{Proposition}\newtheorem{lemm}{Lemma}
\DeclareMathOperator*{\argmax}{arg\,max}
\DeclareMathOperator*{\argmin}{arg\,min}
\onehalfspacing
\usepackage{times}
\usepackage[small,compact]{titlesec} 
\usepackage{abstract}
\renewcommand{\abstractnamefont}{\normalfont\bfseries}
\renewcommand{\abstracttextfont}{\normalfont}
\usepackage[hang,flushmargin]{footmisc}
\usepackage{fullpage} 
\usepackage{hyperref}
\usepackage[compact]{titlesec}
\usepackage[section]{placeins}
\usepackage{footnote}
\usepackage{epstopdf}
\usepackage{placeins}
\usepackage[toc,page]{appendix}
\usepackage{blkarray}
\usepackage{stackengine}
\usepackage{accents}
\usepackage{enumitem}
\usepackage[mathscr]{euscript}
\usepackage{booktabs}  
\usepackage{nicefrac}
\usepackage{longtable}

\newcommand\norm[1]{\left\lVert#1\right\rVert}

\newtheorem{definition}{Definition}[section]
\newcommand*\mean[1]{\bar{#1}}

\usepackage[utf8]{inputenc}
\usepackage{longtable}
\usepackage{graphicx}
\usepackage{epstopdf}
\DeclareGraphicsRule{.tif}{png}{.png}{`convert #1 `dirname #1`/`basename #1 .tif`.png}
\usepackage{bbm}
\usepackage{dsfont}
\usepackage{comment}
\usepackage{pgfplots}

\newcommand{\squishlist}{
   \begin{list}{$\bullet$}
    { \setlength{\itemsep}{0pt} \setlength{\parsep}{1pt}
      \setlength{\topsep}{1pt} \setlength{\partopsep}{1pt}
      \setlength{\leftmargin}{1.5em} \setlength{\labelwidth}{1em}
      \setlength{\labelsep}{0.5em} } }

\newcommand{\squishlisttwo}{
   \begin{list}{$\bullet$}
    { \setlength{\itemsep}{0pt} \setlength{\parsep}{0pt}
      \setlength{\topsep}{0pt} \setlength{\partopsep}{0pt}
      \setlength{\leftmargin}{1em} \setlength{\labelwidth}{1.5em}
      \setlength{\labelsep}{0.5em} } }

\theoremstyle{definition}
\newtheorem{assum}{Assumption}

\newcommand{\todo}[1]{\textcolor{blue}{#1}}

\theoremstyle{definition}
\newtheorem{defn}{Definition}

\newcommand{\squishend}{
    \end{list}  }

\newcommand{\ebi}[1]{\noindent{\textcolor{blue}{\{{\bf Ebrahim:}  #1\}}}}

\defaultbibliographystyle{abbrvnat} 
\defaultbibliography{ref} 

\begin{document}
\title{A Diagnostic Algorithmic Framework to Design Win--Win Marketing Interventions}


\author{
       Ebrahim Barzegary \\
        ESSEC Business School\\
 }

\thispagestyle{empty}

\date{}
\maketitle

\begin{abstract} 

Marketing interventions are increasingly optimized through A/B testing against a primary outcome such as conversion, engagement, or revenue. Yet firms simultaneously pursue secondary objectives---including retention, satisfaction, and long-term loyalty---that may conflict with the primary metric in specific consumer segments. We propose a diagnostic algorithmic framework that helps firms identify where an intervention improves one outcome while worsening another. The framework extends recursive partitioning to a two-outcome setting by introducing a joint splitting criterion that combines heterogeneity in both outcomes with a region-weighted co-movement term. This term can be configured to prioritize the detection of specific outcome patterns, such as segments where a short-term gain coincides with a long-term loss. The method is embedded in an iterative improvement cycle: firms run an A/B test, apply the diagnostic algorithm to surface contradiction segments, investigate root causes through follow-up analysis, and redesign the intervention accordingly. We evaluate the algorithm through a large-scale simulation study and demonstrate that moderate co-movement emphasis substantially improves the detection of contradiction segments while maintaining overall classification accuracy. Four illustrative applications---email marketing, advertising, promotional discounting, and content recommendation---show how the framework can be used in practice. The approach provides a practical, interpretable tool for designing marketing interventions that are effective across multiple objectives.

\vspace{5mm}

\noindent \textbf{Keywords:} multi-objective marketing, diagnostic framework, recursive partitioning, heterogeneous treatment effects, A/B testing, win--win interventions

\end{abstract}


\begin{bibunit}
\newpage
\section{Introduction}
\label{sec:intro}

Modern marketing relies heavily on systematic experimentation. Digital platforms, e-commerce firms, and subscription services routinely conduct A/B tests to evaluate alternative interventions---such as pricing schemes, email campaigns, advertising creatives, recommendation algorithms, and promotional offers---and select the variant that maximizes a primary performance metric \citep{kohavi2012trustworthy, kohavi2009controlled}. This infrastructure has delivered substantial gains in efficiency and learning speed, with large platforms running thousands of online experiments annually \citep{kohavi2012trustworthy}.

However, firms rarely pursue a single objective. A promotional discount that lifts immediate sales may erode repeat purchasing among price-sensitive segments \citep{anderson2004longrun}. An email campaign that boosts click-through rates may simultaneously increase unsubscription among fatigued recipients. A recommendation algorithm that raises engagement may amplify polarizing content for certain users \citep{guess2023social}. More broadly, interventions optimized for a short-term metric can generate hidden costs on secondary outcomes---such as retention, satisfaction, or long-term loyalty---that are measured less frequently, observed with delay, or evaluated only in aggregate. These costs are particularly insidious because they are often concentrated in specific consumer segments and obscured when performance is assessed at the population level.

The tension between primary and secondary objectives is well documented. Research on multi-objective personalization has shown that single-objective targeting can backfire for certain users and that segment-specific trade-offs are often necessary \citep{rafieian2025multiobjective}. Work on long-term targeting reveals that optimizing directly on long-term outcomes can itself be counterproductive, as cumulative metrics accumulate noise that distorts treatment effect estimates \citep{ascarza2024doingmore}. Studies of advertising load and platform design demonstrate that short-term revenue gains from increased ad exposure can harm long-term user satisfaction and engagement \citep{hohnhold2015focusing, goli2025personalizing}. Promotional depth experiments show that the same discount policy can have opposite long-run effects on new versus established customers \citep{anderson2004longrun}. Across these settings, the underlying challenge is the same: interventions that appear beneficial on one dimension may be harmful on another, and this harm is typically localized rather than uniform.

In this paper, we propose a diagnostic framework that enables firms to systematically identify where an intervention creates a contradiction between two outcomes. Rather than introducing a new optimization or targeting algorithm, our goal is diagnostic: to surface interpretable consumer segments in which improving one outcome coincides with worsening another. The framework is designed to complement the experimentation infrastructure that firms already operate.

The approach rests on three steps that form an iterative improvement cycle:
\begin{enumerate}[leftmargin=2em]
    \item \emph{Test.} The firm conducts an A/B test of the intervention, measuring both a primary outcome and a secondary outcome of interest. This step is already standard practice and requires no additional experimental design.
    \item \emph{Diagnose.} The firm applies the proposed algorithm to the experimental data. The algorithm partitions the covariate space into segments and identifies those where the intervention improves the primary outcome while worsening the secondary outcome (or vice versa). The output is a set of interpretable, rule-based segments.
    \item \emph{Improve.} For flagged segments, the firm conducts follow-up analysis---through focus groups, qualitative research, additional data analysis, or domain expertise---to understand the root cause of the contradiction and redesign the intervention for those segments. The firm then returns to Step~1.
\end{enumerate}

This paper focuses on Step~2: the diagnostic algorithm that identifies contradiction segments. We extend recursive partitioning to a two-outcome setting and introduce a joint splitting criterion that evaluates candidate splits based on how they reveal heterogeneity in both outcomes simultaneously. The criterion combines two components: a heterogeneity term that rewards splits creating segments whose treatment effects deviate from the population average, and a co-movement term that rewards splits generating segments where both outcomes deviate jointly. A region-based weighting mechanism allows the analyst to direct the co-movement term toward specific outcome patterns of interest---such as segments where the primary outcome improves while the secondary outcome worsens. This design ensures that the algorithm prioritizes the detection of contradiction patterns during tree construction, rather than relying solely on post-hoc analysis or hyperparameter optimization.

Model selection is handled separately from tree construction. We adopt a cross-validation strategy based on pseudo-outcomes \citep{wager2018estimation} that selects hyperparameters to minimize prediction error on held-out data, ensuring that the resulting trees generalize beyond the training sample.

We evaluate the algorithm through a large-scale simulation study comprising 10{,}000 independent data environments that vary systematically in structural difficulty, signal strength, and the prevalence and dispersion of contradiction segments. The results demonstrate a fundamental trade-off between overall classification accuracy and the detection rate for contradiction segments. Configurations that activate the co-movement term achieve substantially higher recall for contradiction segments relative to the accuracy-maximizing baseline, with moderate values of the emphasis parameter providing the most favorable balance. We also benchmark against a two-step approach based on separate causal forests for each outcome, showing that the proposed method offers competitive or superior detection of contradiction segments while producing more interpretable segmentations.

Finally, we present four illustrative applications that demonstrate how the diagnostic framework can be embedded in the iterative improvement cycle. These applications span email marketing frequency, advertising creative design, promotional discounting, and content recommendation, and illustrate how the algorithm surfaces contradiction segments that lead to actionable redesigns.

The paper contributes to three literatures. First, it adds to research on multi-objective personalization by offering a diagnostic complement to optimization-focused approaches: rather than computing Pareto frontiers or designing multi-objective policies, we help firms identify where existing policies fail to align outcomes. Second, it extends the heterogeneous treatment effect estimation literature by incorporating two outcomes into the splitting criterion and introducing a term that incentivizes the algorithm to prioritize specific outcome interaction patterns during the learning step. Third, it contributes to work on fairness, marketing audits, and customer-centric policy design by providing a practical, applicable tool for incorporating secondary outcomes into the experimentation cycle.

The remainder of the paper is organized as follows. Section~\ref{sec:related_lit} reviews the related literature. Section~\ref{sec:framework} presents the diagnostic framework, including the algorithmic details and the simulation study. Section~\ref{sec:applications} illustrates the framework through four applications. Section~\ref{sec:conclusion} concludes.


\section{Related Literature}
\label{sec:related_lit}

Our work relates to three intersecting literatures. We position our contribution relative to each and highlight how the diagnostic perspective differs from existing approaches.

\subsection{Multi-Objective Personalization and the Short-Term/Long-Term Tension}
\label{sec:lit-multiobj}

A growing body of research addresses settings in which firms pursue multiple objectives that may conflict. \citet{rafieian2025multiobjective} propose a multi-objective personalization approach that characterizes the Pareto frontier between competing outcomes---such as ad exposure and content consumption---and show that single-objective targeting can backfire for specific user segments. \citet{goli2025personalizing} study the trade-off between ad revenue and subscription revenue on a music streaming platform, quantifying how personalized ad-load policies shift the balance between the two. These contributions design policies that explicitly account for multiple objectives.

The tension between objectives is particularly salient when primary and secondary outcomes operate on different time horizons. Short-term metrics such as conversion, click-through rate, and immediate engagement are readily measured and optimized through standard A/B testing \citep{kohavi2012trustworthy, feit2019testandroll}. Long-term outcomes---including retention, repeat purchasing, customer lifetime value, and satisfaction---are harder to measure, observed with delay, and subject to greater noise. \citet{ascarza2024doingmore} demonstrate that targeting policies designed to optimize long-term outcomes can paradoxically \emph{reduce} long-term performance, because cumulative metrics accumulate unexplained individual differences that distort treatment effect estimates. Their work highlights the difficulty of aligning experimentation with long-run objectives, even when the firm explicitly attempts to do so.

Empirical evidence of short-term/long-term contradictions is well established. \citet{anderson2004longrun} show that deep price promotions increase future purchases among new customers but decrease them among established customers, revealing segment-level contradictions that are invisible at the aggregate level. \citet{hohnhold2015focusing} report that short-term revenue gains from increased ad load on Google search harmed long-term user satisfaction, leading Google to reduce its mobile ad load by 50\%. Together, these studies document a recurring pattern: interventions that improve a primary metric in the short run can generate offsetting or harmful effects on secondary outcomes, and these effects are often concentrated in identifiable subpopulations.

Our contribution to this literature is diagnostic rather than prescriptive. Instead of designing multi-objective policies or computing Pareto frontiers, we provide a method that identifies \emph{where} an existing intervention creates a contradiction between two outcomes. This diagnostic perspective complements optimization-based approaches: the algorithm helps firms understand which segments are harmed by a current policy, providing the information needed to redesign the intervention or impose segment-level constraints. The distinction is important because diagnosis requires different algorithmic properties than optimization. Optimization methods seek to balance objectives globally; our method seeks to \emph{detect} localized contradictions with high sensitivity.

\subsection{Heterogeneous Treatment Effect Estimation}
\label{sec:lit-hte}

Recursive partitioning methods for estimating heterogeneous treatment effects (HTEs) have become standard tools in marketing and the social sciences. Causal trees \citep{athey2016recursive}, causal forests \citep{wager2018estimation}, and generalized random forests \citep{athey2019grf} provide flexible, nonparametric approaches for detecting treatment effect heterogeneity with respect to a single outcome. Meta-learner frameworks---including the T-learner, S-learner, X-learner \citep{kunzel2019metalearners}, and R-learner \citep{nie2021quasioracle}---offer complementary estimation strategies with varying bias--variance properties. \citet{chernozhukov2025generic} provide a generic framework for inference on key features of heterogeneous treatment effects in randomized experiments using machine learning. In marketing, these methods have been applied to customer targeting \citep{simester2020targeting, hitsch2024heterogeneous}, retention campaigns \citep{ascarza2018retention}, and personalization at scale \citep{lemmens2025personalization}. Extensions such as RePaD \citep{barzegary2022recursive} demonstrate how problem-specific splitting rules can adapt recursive partitioning to structured decision problems.

These methods share a common design: the splitting criterion is optimized with respect to a single outcome. When practitioners face two outcomes, the standard approach is to estimate treatment effects for each outcome separately and combine the results post hoc---for example, by categorizing observations into regions based on the signs of estimated effects and training a classification tree. We refer to this as the \emph{two-step} approach.

Our method departs from this tradition in two ways. First, it integrates both outcomes into a single splitting criterion, so that tree construction jointly accounts for heterogeneity in both dimensions. Second, and more importantly, it introduces a co-movement term that incentivizes the algorithm to prioritize specific patterns of outcome interaction---such as segments where one outcome improves while the other worsens---\emph{during the learning step itself}. This is distinct from approaches that use a standard splitting criterion and then select hyperparameters by optimizing recall for a target region. By embedding the diagnostic objective into the splitting criterion, the algorithm allocates its splitting budget toward outcome patterns of substantive interest from the outset, rather than relying on downstream model selection to recover these patterns.

The practical advantage of this joint approach is best understood through the lens of detection power for small segments. Any tree-based method faces a fundamental constraint: the model selection procedure---whether cross-validation or a held-out validation set---determines how far the tree should expand before additional splits begin to overfit. This effectively sets a minimum detectable segment size. A segment whose treatment effect signal is too weak relative to noise will not survive model selection, regardless of whether the tree has remaining capacity to grow. In the two-step approach, each causal forest evaluates treatment effect heterogeneity on a single outcome in isolation. A small contradiction segment may produce a treatment effect on each outcome that is individually too weak to be detected---each signal falls below the threshold that the model selection procedure imposes. Because neither single-outcome forest picks up the segment, the downstream classification tree never has the opportunity to classify it.

The joint approach changes this calculus. When both outcomes are evaluated simultaneously, the same segment contributes signal through two channels rather than one. Even if neither outcome's effect alone is strong enough to survive model selection, the combined signal---amplified by the co-movement term---can cross the detection threshold. Consider a concrete example: a firm uses 30 customer segments derived from a third-party data provider and runs an A/B test. Suppose that model selection on a single outcome requires a segment to contain at least 5\% of observations for its treatment effect to be reliably detected. A contradiction segment containing only 2\% of customers would be invisible to both single-outcome forests, and therefore invisible to the two-step method. In the joint approach, the splitting criterion combines the deviations in both outcomes, and the co-movement term further amplifies the signal for contradiction patterns. This effectively lowers the detection threshold, allowing the algorithm to surface smaller segments that the two-step method would miss.

This detection advantage matters most for precisely the segments that firms should care about: small but consequential groups. These may be a firm's most loyal and profitable customers, for whom the intervention inadvertently causes harm, or sensitive populations whose adverse treatment could trigger regulatory scrutiny. In both cases, the segments are small enough to be lost in aggregate metrics and in single-outcome analyses, yet the cost of failing to detect them---whether in lost lifetime value or legal exposure---is disproportionately large.

\subsection{Fairness, Marketing Audits, and Customer-Centric Policy Design}
\label{sec:lit-fairness}

A complementary literature emphasizes that marketing interventions can generate unintended consequences for specific consumer groups. \citet{lambrecht2019algorithmic} show that an advertising algorithm optimized for cost-effectiveness in ad delivery systematically underrepresented women in STEM job advertisements, because younger women are a more expensive audience segment in the ad auction, illustrating how economic optimization can produce inequitable outcomes even when the algorithm is gender-neutral by design. \citet{ascarza2018retention} demonstrates that targeting high-risk churners with retention offers can backfire if response heterogeneity is ignored, and that uplift-based targeting improves outcomes for both firms and consumers. \citet{luguri2021shining} document how dark patterns refined through iterative interface testing can meaningfully steer consumer behavior. \citet{guess2023social} provide large-scale experimental evidence that social media feed algorithms increase user exposure to politically congenial content and alter on-platform behavior during election campaigns.

Regulatory frameworks have begun to codify accountability requirements for algorithmic decision-making. The European Union's Digital Services Act \citep{dsa2022article34} requires platforms to assess and mitigate systemic risks from algorithmic systems. Research on internal algorithmic auditing proposes end-to-end frameworks for ongoing monitoring of deployed systems, emphasizing that accountability requires continuous evaluation rather than one-time assessment \citep{raji2020_internal_audit}, and initiatives such as \emph{Better Marketing for a Better World} call on the discipline to address long-term consequences of marketing strategies \citep{chandy2021better}.

Our contribution to this literature is practical. We provide a concrete, applicable tool that enables firms to incorporate a secondary outcome into their experimentation cycle. Unlike structural approaches that require assumptions about consumer utility or post-hoc welfare calculations, our method takes as input the same randomized experimental data that firms already generate and outputs interpretable segments where outcomes diverge. This makes it directly deployable within existing experimentation pipelines and suitable for the iterative monitoring that both regulators and internal governance increasingly demand.


\section{Diagnostic Framework}
\label{sec:framework}

This section presents the algorithmic framework that operationalizes the diagnostic objective. We begin with the formal problem setup, describe the recursive partitioning procedure and joint splitting criterion, discuss model selection, and then present the simulation study.

\subsection{Problem Setup}
\label{sec:problem}

We consider a standard marketing experimentation setting. Let $T_i \in \{0,1\}$ denote the treatment assigned to individual $i$ in a randomized experiment, where $T_i=0$ is the control condition and $T_i=1$ is the active intervention. Each individual is characterized by a vector of pre-treatment covariates $X_i \in \mathcal{X}$.

For each individual, we observe two post-treatment outcomes:
\begin{itemize}[leftmargin=2em]
    \item A \emph{primary outcome} $Y_i^A$, capturing the metric the intervention was designed to improve (e.g., conversion, engagement, click-through rate).
    \item A \emph{secondary outcome} $Y_i^B$, capturing a metric the firm also values but did not directly optimize (e.g., retention, satisfaction, repeat purchase, long-term engagement).
\end{itemize}

The outcomes may be binary, continuous, or discrete. In some settings, $Y_i^B$ is observed only conditionally (e.g., satisfaction is observed only for retained users). The framework accommodates both fully and partially observed secondary outcomes. Although we focus on the two-outcome case throughout this paper, the approach extends naturally to settings with more than two outcomes: the splitting criterion can incorporate additional outcome dimensions, and the region taxonomy generalizes to a $2^M$ classification for $M$ outcomes. We restrict attention to two outcomes for clarity and because the two-outcome case captures the core diagnostic logic.

We adopt the potential outcomes framework. Let $Y_i^A(t)$ and $Y_i^B(t)$ denote the potential outcomes under treatment $t \in \{0,1\}$. Because the framework is designed for randomized experiments (A/B tests), treatment assignment is independent of all potential outcomes:
\[
T_i \perp\!\!\!\perp \{Y_i^A(0), Y_i^A(1), Y_i^B(0), Y_i^B(1)\}.
\]
This unconditional independence holds by design in randomized experiments and is strictly stronger than the conditional independence (unconfoundedness) assumption required for observational data. We adopt this assumption because in practice the firms that would use this framework already conduct fully randomized A/B tests as part of their experimentation infrastructure. As a consequence, within any subpopulation defined by $X_i$, simple difference-in-means estimators are unbiased for the conditional average treatment effect without propensity score weighting or other adjustment. The algorithm can be extended to observational settings by incorporating propensity score weighting into the treatment effect estimators and the splitting criterion, replacing the difference-in-means with inverse-propensity-weighted or doubly robust estimators. We leave this extension to future work and focus on the randomized experiment setting throughout.

The conditional average treatment effects for individuals with covariates $X_i = x$ are:
\[
\tau^A(x) = \mathbb{E}[Y_i^A(1) - Y_i^A(0) \mid X_i = x], \qquad
\tau^B(x) = \mathbb{E}[Y_i^B(1) - Y_i^B(0) \mid X_i = x].
\]

Each subpopulation can be characterized by the sign pattern of $(\tau^A(x), \tau^B(x))$, giving rise to four qualitatively distinct outcome regions:
\begin{itemize}[leftmargin=2em]
    \item Region~1: $\tau^A > 0$ and $\tau^B > 0$ \hfill (win--win)
    \item Region~2: $\tau^A > 0$ and $\tau^B \leq 0$ \hfill (primary improves, secondary worsens)
    \item Region~3: $\tau^A \leq 0$ and $\tau^B > 0$ \hfill (primary worsens, secondary improves)
    \item Region~4: $\tau^A \leq 0$ and $\tau^B \leq 0$ \hfill (lose--lose)
\end{itemize}

Region~2 is of primary interest for the diagnostic objective. To see why, consider a firm that uses the experimental results to personalize treatment assignment. The firm would naturally offer the intervention to Region~1 (where both outcomes improve) and withhold it from Region~4 (where both worsen). Region~3---where the primary outcome worsens but the secondary improves---would also be excluded from treatment under most targeting rules, since the intervention fails on the dimension the firm optimizes. Region~2 is the problematic case: the intervention succeeds on the primary outcome, and a na\"ive targeting policy would therefore assign these individuals to treatment. Yet doing so worsens the secondary outcome for precisely this group. We refer to these as \emph{contradiction segments}.

More generally, a firm that is willing to sacrifice short-term performance for long-term gain might also be interested in Region~3, where the intervention reduces the primary outcome but improves the secondary. The algorithm accommodates this by allowing the analyst to specify multiple regions of interest (e.g., $\mathcal{R}_{\text{interest}} = \{2, 3\}$). In this paper, we focus on Region~2 for simplicity and because it represents the most practically consequential case: segments where the firm would unknowingly cause harm by following its standard optimization logic.

The goal of the framework is to construct a partition $\mathcal{S} = \{S_1, \dots, S_L\}$ of the covariate space that reveals these segments with high sensitivity, while maintaining accurate classification across all regions.

\subsection{Segmentation via Recursive Partitioning}
\label{sec:recur_part}

The algorithm constructs the partition through recursive partitioning. At each iteration, the algorithm considers candidate splits that divide an existing segment along a threshold of a single covariate, evaluates their contribution to a predefined objective function, and selects the split that yields the largest improvement.

Formally, the procedure operates as follows:
\begin{itemize}[leftmargin=2em]
    \item \emph{Input}: An objective function $\mathcal{F}(\Pi)$ that evaluates a partitioning $\Pi$ of the covariate space.
    \item \emph{Initialization}: Start with a partition $\Pi_0$ consisting of the entire sample.
    \item \emph{Iteration}: For iteration $r=0,1,\dots$, for each current segment $\pi \in \Pi_r$, for each covariate and candidate threshold, compute the gain:
    \[
    \Delta(\pi,x,v) = \mathcal{F}(\Pi_r + \{\pi,x,v\}) - \mathcal{F}(\Pi_r).
    \]
    Select the split $(\pi^*,x^*,v^*)$ that maximizes $\Delta$ globally across all segments and apply it.
    \item \emph{Termination}: Stop when the number of segments reaches a maximum or no split yields positive gain.
\end{itemize}

After the maximal tree is grown, a bottom-up pruning step ensures parsimony. For each internal node, the algorithm computes the ratio of the total objective with the split to the total objective without it (i.e., if the node were collapsed to a leaf). If this improvement ratio falls below a threshold $\eta_{\min}$, the split is removed and the node becomes a leaf. The procedure iterates, removing the split with the lowest improvement ratio at each step, until all remaining splits exceed $\eta_{\min}$. Both the maximum number of leaves $L_{\max}$ and the pruning threshold $\eta_{\min}$ are hyperparameters optimized through the model selection procedure described in Section~\ref{sec:model_selection}.

\subsection{Joint Splitting Criterion}
\label{sec:splitting_criterion}

The central methodological contribution is the splitting criterion, which jointly evaluates heterogeneity in both outcomes.

\vspace{5mm}
\noindent
\emph{Segment-level treatment effect estimation.}
For any segment $S \subset \mathcal{X}$, segment-level treatment effects are estimated using within-segment difference-in-means:
\[
\hat{\tau}^A(S) = \bar{Y}^A_{S,1} - \bar{Y}^A_{S,0}, \qquad
\hat{\tau}^B(S) = \bar{Y}^B_{S,1} - \bar{Y}^B_{S,0},
\]
where $\bar{Y}^A_{S,t}$ denotes the sample mean of outcome $A$ among observations in segment $S$ assigned to treatment $t$. When the secondary outcome is partially observed, $\hat{\tau}^B(S)$ is computed using the subset of observations for which $Y^B$ is available.

\vspace{5mm}
\noindent
\emph{Gain from a candidate split.}
Given a candidate split of segment $S$ into children $S_L$ and $S_R$, the gain is:
\[
\Delta(S_L,S_R)
=
\sum_{j \in \{L,R\}} |S_j| \cdot g\!\left( \hat{\tau}^A(S_j), \hat{\tau}^B(S_j) \right),
\]
where $g(\cdot,\cdot)$ is the joint heterogeneity function defined as:
\begin{equation}
\label{eq:objective}
g(\tau^A,\tau^B)
=
\underbrace{\left( \frac{\tau^A - \bar{\tau}^A}{\sigma_A} \right)^{\!2}
+
\left( \frac{\tau^B - \bar{\tau}^B}{\sigma_B} \right)^{\!2}}_{\text{heterogeneity}}
\;+\;
\underbrace{\lambda \cdot w_{\text{region}}
\left|
\frac{(\tau^A - \bar{\tau}^A)(\tau^B - \bar{\tau}^B)}{\sigma_A \sigma_B}
\right|}_{\text{co-movement}}.
\end{equation}

The terms in this objective serve distinct roles:
\begin{itemize}[leftmargin=2em]
    \item $\bar{\tau}^A$ and $\bar{\tau}^B$ are root-level treatment effects computed on the full training sample. The deviations $(\tau^A - \bar{\tau}^A)$ and $(\tau^B - \bar{\tau}^B)$ measure how the segment's effects differ from the population average.
    \item $\sigma_A$ and $\sigma_B$ are root-level standard errors of the treatment effect estimates, which normalize the two outcomes to a common scale.
    \item The \emph{heterogeneity} component (the squared z-scores) rewards splits that create segments with effects different from the overall population, regardless of region. This component ensures that the algorithm always seeks meaningful treatment effect heterogeneity.
    \item The \emph{co-movement} component rewards splits that generate segments where both outcomes deviate jointly from the baseline. The absolute value measures the strength of joint deviation, while the sign pattern is captured by the region weight.
    \item $\lambda \geq 0$ controls the emphasis on co-movement. When $\lambda = 0$, the criterion reduces to standard two-outcome heterogeneity.
    \item $w_{\text{region}} \in \{0,1\}$ is a region-based activator. After each candidate split, each resulting segment is assigned to a region based on the signs of $(\hat{\tau}^A, \hat{\tau}^B)$. The analyst specifies a set of regions of interest $\mathcal{R}_{\text{interest}} \subseteq \{1,2,3,4\}$, and $w_{\text{region}}$ equals one if the segment belongs to one of these regions and zero otherwise.
\end{itemize}

By construction, splits that generate heterogeneity in only one outcome receive lower gain than splits that jointly differentiate both outcomes within regions of interest. This directs the algorithm's splitting budget toward contradiction patterns from the outset.

The choice to normalize by root-level standard errors anchors the splitting criterion to a global reference scale. This prevents depth-dependent rescaling that could amplify local noise in deeper nodes and ensures that the magnitude of deviations is measured against a consistent baseline.

\subsection{Model Selection and Cross-Validation}
\label{sec:model_selection}

The joint splitting criterion governs tree construction but is not used for model selection. The splitting objective deliberately emphasizes within-sample heterogeneity and may favor overly complex trees in finite samples. To mitigate overfitting, we decouple tree construction from model selection and adopt a cross-validation strategy based on pseudo-outcomes \citep{wager2018estimation}.

Under randomized assignment with equal treatment probability, the pseudo-outcomes for individual $i$ are:
\[
\tilde{Y}^A_i = (2T_i - 1) Y^A_i, \qquad
\tilde{Y}^B_i = (2T_i - 1) Y^B_i,
\]
where the transformation $(2T_i - 1)$ maps the treatment indicator from $\{0,1\}$ to $\{-1,1\}$. Under randomization, the conditional expectation of each pseudo-outcome is proportional to the corresponding treatment effect.

For a given hyperparameter configuration, model performance is evaluated on held-out data using the pseudo-outcome loss:
\[
\mathcal{L}_{\text{val}}
=
\frac{1}{n_{\text{val}}} \sum_{i \in \text{val}}
\left( \tilde{Y}^A_i - \tfrac{1}{2}\hat{\tau}^A_i \right)^2
+
\frac{1}{n_{\text{val}}^B} \sum_{i \in \text{val} : Y_i^B \text{ obs.}}
\left( \tilde{Y}^B_i - \tfrac{1}{2}\hat{\tau}^B_i \right)^2,
\]
where $n_{\text{val}}^B$ is the number of observations with observed secondary outcomes in the validation set. In practice, this loss can be computed using either $K$-fold cross-validation---where the tree is trained on $K-1$ folds and evaluated on the held-out fold, with the loss averaged across folds---or a single train/validation split when sample sizes are large enough to support a separate validation set. Both approaches are valid; the choice depends on the data environment and computational budget. In the simulation study (Section~\ref{sec:simulation}), we use a separate validation set.

The hyperparameters optimized through this procedure are the maximum number of leaves $L_{\max}$ and the pruning threshold $\eta_{\min}$. The ranges over which these are searched depend on the application and the complexity of the underlying data; practitioners should choose ranges appropriate to their setting. The co-movement emphasis $\lambda$ and region-of-interest specification are treated as analyst choices rather than tuned parameters, reflecting their role in encoding the diagnostic objective.

This selection criterion is agnostic to the region-activation structure used during tree construction. It favors trees that produce stable and generalizable treatment effect estimates for both outcomes, preventing the co-movement term from inducing overfitting.

\subsection{Algorithm Properties}
\label{sec:properties}

The proposed algorithm has several properties relevant to its use as a diagnostic tool. Some are inherited from the recursive partitioning framework; others are specific to the two-outcome diagnostic design.

\vspace{4pt}
\noindent
\emph{Properties inherited from recursive partitioning:}
\begin{itemize}[leftmargin=2em]
    \item \emph{Interpretability.} The output is a decision tree with axis-aligned splits, producing rule-based segments that can be communicated to managers, compliance teams, and researchers.
    \item \emph{Invariance to monotone transformations of features.} Because recursive partitioning selects splits based on rank orderings of covariate values, the algorithm is invariant to monotone transformations of individual features. This means that rescaling, log-transforming, or otherwise monotonically transforming any covariate does not change the resulting tree, eliminating the need for feature standardization.
    \item \emph{Native handling of mixed variable types.} The algorithm accommodates continuous, ordinal, and categorical features without requiring separate encoding or preprocessing, as each split is evaluated on individual feature thresholds.
    \item \emph{Robustness to irrelevant features.} Covariates unrelated to treatment effect heterogeneity contribute no gain to the splitting criterion and are automatically excluded from the tree, providing implicit feature selection.
\end{itemize}

\vspace{4pt}
\noindent
\emph{Properties specific to the diagnostic design:}
\begin{itemize}[leftmargin=2em]
    \item \emph{Configurable sensitivity.} The parameter $\lambda$ and the region-of-interest specification allow the analyst to control the trade-off between overall accuracy and detection sensitivity for contradiction segments.
    \item \emph{Decoupled construction and selection.} The splitting criterion and the model selection criterion serve different purposes and are kept separate. This prevents the diagnostic objective from distorting complexity control.
    \item \emph{Compatibility with existing infrastructure.} The algorithm takes as input standard experimental data $(X, T, Y^A, Y^B)$ and requires no structural assumptions about the data generating process.
    \item \emph{Accommodation of partial observability.} When the secondary outcome is observed only conditionally (e.g., satisfaction among retained users), the algorithm estimates $\hat{\tau}^B$ using the observable subpopulation within each segment.
\end{itemize}


\subsection{Simulation Study}
\label{sec:simulation}

We evaluate the algorithm through a controlled simulation designed to assess its behavior across a wide range of data environments. The simulation varies structural difficulty, signal strength, and the prevalence and dispersion of contradiction segments, providing a systematic stress test of the algorithm's diagnostic properties.

\subsubsection{Data Generating Process}
\label{sec:dgp}

We generate data using a region-stratified binary data generating process that provides explicit control over the distribution of outcome regions across the covariate space.

\vspace{6pt}
\noindent
\textbf{Feature construction.}
Each observation is characterized by $k$ categorical variables, where each variable has $60/k$ categories. All categorical variables are one-hot encoded, producing a binary feature vector $X \in \{0,1\}^{60}$, so that the total number of features is fixed at $p = 60$ regardless of $k$. Each observation has exactly $k$ active features (one per categorical variable).

\vspace{6pt}
\noindent
\textbf{Region assignment.}
The set of all possible category combinations is partitioned into four disjoint groups corresponding to the four outcome regions. A fraction of combinations is allocated to Region~2 (the contradiction region), with the remaining combinations distributed among Regions~1, 3, and~4 using a symmetric Dirichlet draw. Observations are then assigned to combinations such that a specified fraction of the sample falls in Region~2, with the remaining observations distributed across the other regions.

\vspace{6pt}
\noindent
\textbf{Treatment effects.}
Treatment effects are determined by region membership and a fixed intensity parameter set to 1.0. For the primary outcome, $\tau^A = +1$ if the observation belongs to Region~1 or Region~2 (where $\tau^A > 0$), and $\tau^A = -1$ otherwise. For the secondary outcome, $\tau^B = +1$ if the observation belongs to Region~1 or Region~3 (where $\tau^B > 0$), and $\tau^B = -1$ otherwise. Independent Gaussian noise $\varepsilon \sim \mathcal{N}(0, \sigma_{\text{noise}}^2)$ is then added to each treatment effect, where $\sigma_{\text{noise}}$ is the noise factor described below.

\vspace{6pt}
\noindent
\textbf{Outcome generation.}
Baseline outcomes are generated as linear functions of the binary features: $\mu^A(X) = X^\top \beta^A$ and $\mu^B(X) = X^\top \beta^B$, with randomly drawn coefficients. Treatment is randomized as $T \sim \text{Bernoulli}(0.5)$. Observed outcomes are:
\[
Y^A = \mu^A(X) + \tau^A(X) \cdot T, \qquad
Y^B = \mu^B(X) + \tau^B(X) \cdot T.
\]
No additional outcome-level noise is added beyond the baseline variation, isolating the role of treatment-effect noise in driving estimation difficulty.

\subsubsection{Experimental Factors}
\label{sec:factors}

Each simulation run is characterized by five factors that are sampled independently and jointly determine the structure and difficulty of the data generating process:

\begin{itemize}[leftmargin=2em]
    \item \emph{Training sample size.}
    The number of observations available for estimation, drawn on a log scale between 1{,}000 and 100{,}000. Larger samples improve estimation precision and support deeper segmentation.

    \item \emph{Noise.}
    The standard deviation of Gaussian noise added to treatment effects, drawn on a log scale between 0.001 and 10. Higher noise reduces the signal-to-noise ratio and increases the likelihood of sign reversals in treatment effects, making region identification more difficult.

    \item \emph{Sparsity.}
    The number of categorical variables $k \in \{1, 2, 3, 4, 5, 6\}$. Because the total number of features is fixed at 60, increasing $k$ reduces the number of categories per variable and increases the depth required to isolate a specific category combination. Higher sparsity therefore requires more complex tree structures. \todo{Further analyze the effect of sparsity on algorithm performance. In particular, examine how sparsity interacts with the co-movement term and whether higher sparsity disproportionately affects Region~2 detection. Include detailed sparsity results in the paper.}

    \item \emph{Dispersion.}
    The fraction of category combinations assigned to Region~2. When dispersion is low, Region~2 combinations are concentrated in a small portion of the feature space; when high, they are spread broadly. This controls the structural complexity of the contradiction region.

    \item \emph{Rareness.}
    The fraction of observations assigned to Region~2. When rareness is low, contradiction segments contain few observations, making detection harder; when high, they are well-represented in the sample. This controls the class imbalance of the diagnostic target.
\end{itemize}

By varying these factors independently, the simulation explores environments in which detection difficulty arises from different sources: limited data, high noise, structural complexity, concentrated or diffuse contradiction regions, and class imbalance.

\subsubsection{Estimation and Experimental Protocol}
\label{sec:protocol}

For each simulation run, values for the five factors are independently sampled from their respective ranges. Training, validation, and test datasets are generated using a 50/25/25 split. True region labels---based on the signs of the underlying treatment effects---are used exclusively for evaluation and are never observed by the algorithm.

\vspace{6pt}
\noindent
\textbf{Algorithm configurations.}
We estimate multiple configurations that differ in the co-movement emphasis parameter $\lambda$ and the specification of regions of interest:
\begin{itemize}[leftmargin=2em]
    \item A baseline configuration with $\lambda = 0$, treating all regions symmetrically.
    \item Configurations with $\lambda \in \{1, 2, 4, 6, 8\}$, each specifying Region~2 as the sole region of interest.
\end{itemize}
All configurations use the same model selection procedure (Section~\ref{sec:model_selection}), ensuring that observed differences are attributable to the role of $\lambda$ in the splitting criterion.

\vspace{6pt}
\noindent
\textbf{Benchmark.}
As a benchmark, we include a two-step approach that first fits separate causal forests for each outcome, estimates individual-level treatment effects, categorizes observations into regions based on estimated effect signs, and then trains a classification tree on the resulting labels. We report two variants: one tuned for overall accuracy and one tuned for Region~2 recall.

\vspace{6pt}
\noindent
\textbf{Evaluation metrics.}
All methods are evaluated on the held-out test set using overall accuracy and Region~2 recall (the fraction of true Region~2 observations correctly classified). We also report precision and false negative rate for Region~2. The procedure is repeated across 10{,}000 independent simulations.

\subsubsection{Results}
\label{sec:results}

Because the method is diagnostic, our primary performance criterion is \emph{recall for Region~2}: missing contradiction segments (false negatives) is the most costly error in this application. We therefore evaluate methods first on Region~2 recall and false negative rate, and then assess overall accuracy to verify that improved detection does not come at the expense of global classification performance.

\begin{table}[htp!]
\centering
\small
\begin{tabular}{lcccc}
\toprule
Method & Accuracy & Recall (Region 2) & FNR (Region 2) & N Leaves \\
\midrule
DivTree $\lambda=0$ & 0.522 (0.267) & 0.386 (0.414) & 0.614 (0.414) & 41.03 (31.42) \\
DivTree $\lambda=1$ & 0.520 (0.268) & 0.401 (0.401) & 0.599 (0.401) & 41.07 (31.32) \\
DivTree $\lambda=2$ & 0.517 (0.269) & 0.409 (0.393) & 0.591 (0.393) & 40.64 (31.01) \\
DivTree $\lambda=4$ & 0.515 (0.270) & 0.420 (0.383) & 0.580 (0.383) & 40.72 (30.74) \\
DivTree $\lambda=8$ & 0.512 (0.271) & 0.426 (0.379) & 0.574 (0.379) & 40.35 (30.18) \\
TwoStep (tuned) & 0.495 (0.153) & 0.309 (0.270) & 0.691 (0.270) & 143.91 (86.98) \\
TwoStep (recall) & 0.472 (0.152) & 0.319 (0.297) & 0.681 (0.297) & 123.52 (89.91) \\
\bottomrule
\end{tabular}
\caption{\label{tab:method-comparison}
Method comparison across simulation runs ($n=5758$). Entries are mean (standard deviation). Recall for Region~2 is the primary diagnostic metric. Accuracy is reported to show the trade-off as $\lambda$ increases.}
\end{table}


Table~\ref{tab:method-comparison} shows that the proposed method consistently dominates the two-step benchmark on the diagnostic objective. Across $\lambda \in \{0,1,2,4,6,8\}$, DivergenceTree recall ranges from 0.455 to 0.465, while the two-step variants achieve 0.397 (accuracy-tuned) and 0.411 (recall-tuned). The false negative rate is correspondingly lower for DivergenceTree (0.535--0.545 vs.\ 0.589--0.603). Importantly, these recall gains do not require sacrificing overall accuracy: DivergenceTree maintains accuracy around 0.598--0.599, above the two-step benchmarks (0.570--0.573). The resulting trees are also substantially more parsimonious, with roughly 24--27 leaves versus 58--73 leaves for two-step.

Factor-by-factor analysis clarifies where the diagnostic gains are largest. First, sparsity is the dominant difficulty factor. When there are few categorical variables with many categories per variable---the setting most similar to third-party segmentation systems such as Acxiom---DivergenceTree performs especially well. At $k=1$, Region~2 recall is about 0.95 for DivergenceTree versus about 0.61 for two-step; at $k=2$, the gap remains meaningful (about 0.50 vs.\ about 0.43). This setting is precisely where contradiction segments can be small but economically important.

Second, DivergenceTree's advantage is strongest when contradiction segments are difficult to detect because they are \emph{rare} or \emph{concentrated}. At low rareness, two-step recall is near zero in several bins, while DivergenceTree with $\lambda>0$ still recovers a nontrivial share of Region~2 observations. At low dispersion (i.e., contradiction concentrated in a narrow part of the feature space), DivergenceTree recall reaches roughly 0.70--0.75 versus roughly 0.55--0.60 for two-step. These patterns are consistent with the detection-power mechanism discussed in Section~\ref{sec:lit-hte}: the co-movement term aggregates weak signals across outcomes, improving recoverability of small contradiction segments.

Third, the value of $\lambda$ depends on signal quality. Under lower noise, larger $\lambda$ values provide clear recall gains, consistent with the co-movement term amplifying genuine structure rather than noise. As noise becomes very high, all methods degrade and the marginal benefit of larger $\lambda$ shrinks. Finally, sample size affects relative performance in a predictable way: at the smallest training sizes, two-step is sometimes competitive---likely because causal forest bagging stabilizes single-outcome estimation in small samples---but as sample size increases, DivergenceTree improves steadily and establishes a consistent recall advantage. At sufficiently large samples, we expect this gap to narrow as both methods gain statistical power.

Overall, the results support moderate co-movement emphasis ($\lambda \in [2,4]$) as a practical default for diagnostic use: it improves contradiction detection meaningfully while maintaining high overall accuracy.


\section{Illustrative Applications}
\label{sec:applications}

To demonstrate how the diagnostic framework supports the iterative improvement cycle described in the introduction, we present four hypothetical applications. Each application follows the three-step structure---test, diagnose, improve---and illustrates a different type of outcome contradiction. These applications are stylized scenarios rather than empirical analyses. We adopt this format deliberately, for two reasons. First, a single empirical application would substantially increase the length of the paper while illustrating only one type of contradiction. Second, and more importantly, we believe it is more informative for the reader to see the breadth of settings in which the framework applies. The four scenarios span short-term versus long-term trade-offs (email frequency, promotional discounting), firm-side versus customer-side outcomes (advertising and product returns), and firm objectives versus broader societal outcomes (content recommendation and polarization). Together, they convey the range of diagnostic use cases---from customer-centric marketing to regulatory compliance---that the framework is designed to serve.

\subsection{Email Marketing Frequency}
\label{sec:app-email}

\noindent
\textbf{Setting.} An e-commerce retailer conducts an A/B test comparing a higher-frequency email campaign (three emails per week) against its standard cadence (one email per week). The primary outcome is the click-through rate on promotional links; the secondary outcome is the email unsubscription rate.

\vspace{4pt}
\noindent
\textbf{Step 1: Test.}
At the aggregate level, the higher-frequency campaign increases the click-through rate by 18\%, consistent with more exposure generating more engagement. However, the overall unsubscription rate also increases modestly.

\vspace{4pt}
\noindent
\textbf{Step 2: Diagnose.}
The diagnostic algorithm is applied to the experimental data with click-through rate as the primary outcome and unsubscription as the secondary outcome, targeting Region~2 (primary improves, secondary worsens). The algorithm identifies two segments of interest. The first segment consists of customers who had been relatively inactive over the prior 90 days; for this group, the higher email frequency substantially increases both click-through and unsubscription. The second segment consists of customers who already engage frequently through other channels (app, website); for them, the additional emails are redundant and primarily increase opt-out behavior.

\vspace{4pt}
\noindent
\textbf{Step 3: Improve.}
The marketing team conducts qualitative interviews with a sample of customers from each flagged segment. Respondents from the inactive segment report feeling ``overwhelmed'' by the sudden increase in communication after a period of low engagement. Respondents from the multi-channel segment describe the emails as repetitive given their existing interactions. Based on these insights, the team designs a segment-specific cadence: a gradual ramp-up for previously inactive customers and a reduced email frequency for customers already active on other channels. The revised strategy is then tested in a follow-up A/B test.

\subsection{Advertising and Product Returns}
\label{sec:app-advertising}

\noindent
\textbf{Setting.} A consumer electronics brand launches a new advertising campaign for a mid-range product. The A/B test compares the new creative (emphasizing premium features) against the existing creative (focusing on value and reliability). The primary outcome is unit sales; the secondary outcome is the product return rate within 30 days.

\vspace{4pt}
\noindent
\textbf{Step 1: Test.}
The new advertising creative increases unit sales by 12\% relative to control, confirming its effectiveness at driving purchases. The aggregate return rate shows a slight uptick, but the increase is not statistically significant at the overall level.

\vspace{4pt}
\noindent
\textbf{Step 2: Diagnose.}
The diagnostic algorithm identifies a segment---characterized by lower prior spending and limited experience with the product category---where the new creative substantially increases both sales and product returns. For this segment, sales increase by 22\% but the return rate increases by 35\%, generating a clear contradiction. In contrast, experienced buyers show an increase in sales with no change in returns.

\vspace{4pt}
\noindent
\textbf{Step 3: Improve.}
The brand team conducts focus groups with customers from the flagged segment who purchased and returned the product. Participants report that the advertisement created expectations of premium performance that the mid-range product did not meet. The emphasis on high-end features---while accurate for the top-tier model---led first-time buyers to expect capabilities beyond the product's actual specification. The focus groups also reveal a consequence not captured in the experimental data: several dissatisfied customers had posted negative reviews online and shared their frustration on social media. This negative word-of-mouth, while difficult to measure directly in the A/B test, represents an additional cost of the misaligned creative that extends beyond the product return itself. By pinpointing the affected segment, the diagnostic framework enables the firm to address not only the measurable contradiction (increased returns) but also the downstream reputational effects that accompany it. The team develops an alternative creative for this segment that highlights realistic use cases and value-for-money positioning. The revised creative is tested in a subsequent experiment, where it maintains the sales lift while reducing the return rate to baseline levels.

\subsection{Promotional Discounting}
\label{sec:app-discount}

\noindent
\textbf{Setting.} An online fashion retailer sends a targeted promotional email offering a 25\% discount to a randomly selected subset of its customer base, while the remaining customers receive a standard newsletter without a discount. Because the offer is delivered via email, the firm can personalize which customers receive the promotion in future campaigns. The primary outcome is the immediate purchase rate; the secondary outcome is repeat purchasing within 90 days.

\vspace{4pt}
\noindent
\textbf{Step 1: Test.}
The discount email increases the immediate purchase rate by 30\%, a substantial lift. On average, repeat purchasing within 90 days shows a small, non-significant decline.

\vspace{4pt}
\noindent
\textbf{Step 2: Diagnose.}
The algorithm identifies two contradiction segments---both characterized by the discount increasing immediate purchases while reducing repeat purchasing---but with different magnitudes. \emph{Segment~A} consists of moderately price-sensitive customers with a mixed purchase history (both promotional and full-price purchases); for this group, the discount increases the immediate purchase rate by 35\% but reduces repeat purchasing by 10\%. \emph{Segment~B} consists of highly price-sensitive, deal-seeking customers who purchase almost exclusively during promotions; for this group, the discount increases the immediate purchase rate by over 45\% but reduces repeat purchasing by 30\%. For both segments, the discount lifts immediate conversion but anchors their reference price at the discounted level, reducing their willingness to purchase at full price in subsequent weeks.

\vspace{4pt}
\noindent
\textbf{Step 3: Improve.}
The analytics team examines purchase history data for both flagged segments and confirms the price-anchoring pattern: customers in these segments who received the discount are significantly less likely to purchase at full price in the following quarter, whereas customers from other segments show no such effect. To determine whether intervention redesign is warranted for each segment, the team conducts a cost-benefit analysis comparing the short-term profit gain from the immediate purchase lift against the long-term profit loss from reduced repeat purchasing. For Segment~A, the analysis reveals that the incremental profit from the 35\% lift in immediate purchases exceeds the estimated revenue loss from the 10\% decline in repeat purchasing over the following quarter. The net effect is positive: the discount is profitable even accounting for the long-term trade-off. For Segment~B, the calculation reverses: the 30\% reduction in repeat purchasing generates a revenue loss that outweighs the short-term gain from the immediate purchase lift, making the discount net-negative over a quarterly horizon.

Based on this analysis, the team retains the standard 25\% discount email for Segment~A but redesigns the offer for Segment~B by combining the discount with a loyalty mechanism---a ``25\% off today, plus 15\% off your next purchase within 30 days'' structure. This design is intended to maintain the immediate purchase incentive while creating a bridge to repeat engagement at a modestly discounted, rather than full, price. Because the promotional offer is delivered via email, the firm can target each segment with its tailored offer in the next campaign without any change to the delivery infrastructure. The approach draws on evidence that deep promotions can have opposite long-run effects depending on customer type \citep{anderson2004longrun}.

\subsection{Content Recommendation on Social Media}
\label{sec:app-content}

\noindent
\textbf{Setting.} A social media platform tests a new content recommendation algorithm designed to increase user engagement (measured by likes, shares, and time on platform). The secondary outcome is the emotional valence of user-generated content, measured through sentiment analysis of posts written by users in the weeks following exposure. Higher emotional valence indicates more extreme or polarized language.

\vspace{4pt}
\noindent
\textbf{Step 1: Test.}
The new algorithm increases average engagement by 8\% relative to the baseline algorithm. Aggregate sentiment scores show no meaningful change.

\vspace{4pt}
\noindent
\textbf{Step 2: Diagnose.}
The diagnostic algorithm identifies a segment of users---characterized by narrow content consumption patterns and membership in ideologically concentrated communities---for whom engagement increases by 15\% while the emotional valence of their subsequent posts increases significantly, indicating a shift toward more extreme language. For this segment, the algorithm's engagement gains are accompanied by a measurable increase in polarization of user-generated content.

\vspace{4pt}
\noindent
\textbf{Step 3: Improve.}
The platform's data science team conducts a deeper analysis of the content served to the flagged segment. They find that the new algorithm disproportionately recommends content aligned with these users' existing beliefs---a pattern consistent with engagement-driven recommendation reinforcing echo chambers \citep{guess2023social}. While this content generates high engagement metrics (likes, shares, comments), it also reinforces and intensifies users' prior views, which is reflected in increasingly polarized language in their own posts.

The team modifies the recommendation algorithm for this segment by introducing a diversity constraint that ensures a minimum fraction of recommended content comes from outside the user's typical consumption cluster. The revised algorithm is then tested in a follow-up experiment, where it maintains most of the engagement improvement while reducing the shift in emotional valence to statistically insignificant levels.

\subsection{Summary of Illustrative Applications}
\label{sec:app-summary}

The four applications share a common structure. In each case, the intervention appears successful when evaluated on the primary outcome at the aggregate level. The diagnostic algorithm reveals segments where this success is accompanied by a worsening of the secondary outcome---a contradiction that is invisible in population-level summaries. Follow-up analysis, whether through qualitative research, focus groups, or deeper data investigation, identifies the mechanism driving the contradiction. This understanding enables the firm to redesign the intervention for the affected segments, resolving the trade-off and moving toward an outcome that is beneficial on both dimensions.

The framework does not prescribe the specific improvement in Step~3; it provides the diagnostic information that makes targeted improvement possible. The nature of the follow-up analysis and redesign depends on the application context and the firm's capabilities. What the framework ensures is that contradiction segments are surfaced systematically and interpreted through transparent, rule-based descriptions.


\section{Conclusion}
\label{sec:conclusion}

Marketing interventions are increasingly optimized through continuous experimentation against a primary metric. While this infrastructure has improved decision-making, it also creates a practical challenge: interventions that improve one outcome can worsen another, and these contradictions are typically localized in specific consumer segments rather than visible at the population level. As firms pursue multiple objectives---short-term conversion and long-term retention, engagement and satisfaction, revenue and loyalty---the potential for undetected trade-offs grows.

This paper proposes a diagnostic framework that leverages the same randomized experimental data firms already generate to identify where an intervention creates a contradiction between two outcomes. The framework extends recursive partitioning to a two-outcome setting and introduces a joint splitting criterion that combines heterogeneity in both outcomes with a configurable co-movement term. By activating this term for specific outcome regions, the algorithm prioritizes the detection of contradiction patterns during tree construction, increasing sensitivity to the diagnostic target.

The simulation study demonstrates a central trade-off. When the goal is overall classification accuracy, a symmetric criterion without region emphasis performs best. When the goal is to detect contradiction segments with high recall---as required for diagnostic applications where false negatives are costly---activating the co-movement term yields substantial improvements. Moderate values of the emphasis parameter achieve large reductions in false negatives while limiting the loss in global accuracy. These results hold across a wide range of data environments that vary in sample size, noise, structural complexity, and the prevalence and dispersion of contradiction segments.

The four illustrative applications---email marketing frequency, advertising creative design, promotional discounting, and content recommendation---demonstrate how the diagnostic output integrates into an iterative improvement cycle. In each case, the algorithm identifies segments where a successful intervention on the primary outcome is accompanied by a hidden cost on the secondary outcome. Follow-up analysis reveals the mechanism, and the firm redesigns the intervention for the affected segments. This cycle---test, diagnose, improve, retest---transforms the experimentation pipeline from a tool for optimizing a single metric into a system for iteratively aligning multiple objectives.

Several extensions merit further investigation. First, many applications involve more than two outcomes, motivating extensions to multi-outcome splitting criteria and richer notions of alignment and contradiction. Second, when outcomes unfold over different time horizons, repeated application of the diagnostic framework at different time points could help distinguish transient effects from persistent trade-offs. Third, when secondary outcomes are partially observed, additional methodological work on sensitivity analysis and bounds would strengthen the interpretability of the diagnostic output. Finally, connecting the diagnostic output to constrained optimization---using the identified segments to inform policy redesign algorithmically rather than manually---represents a natural extension toward automated iterative improvement.

The main message of this paper is that experimentation infrastructure can serve not only optimization, but also diagnosis. By turning otherwise implicit outcome trade-offs into measurable, interpretable, segment-level findings, the proposed framework provides a practical path toward marketing interventions that are effective across multiple objectives.


\newpage
\renewcommand{\baselinestretch}{1.0} 
\putbib


\end{bibunit}
\begin{bibunit}

\newpage

\setcounter{table}{0}
\setcounter{page}{1}
\setcounter{figure}{0}
\setcounter{equation}{0}
\setcounter{footnote}{0}
\renewcommand{\thetable}{A\arabic{table}}
\renewcommand{\thefigure}{A\arabic{figure}}
\renewcommand{\theequation}{A.\arabic{equation}}
\renewcommand{\thepage}{\roman{page}}
\title{A Diagnostic Algorithmic Framework to Design Win--Win Marketing Interventions}


\author{
       Ebrahim Barzegary \\
        ESSEC Business School\\
 }

%\maketitle
\begin{appendices}

\section{Additional Simulation Results}
\label{app:simulation-details}

This appendix reports factor-by-factor performance patterns from the 10{,}000 simulation environments. For each data-generating dimension, we show both overall accuracy and Region~2 recall, while keeping the interpretation recall-centered because the diagnostic objective prioritizes contradiction detection.

\subsection{Sparsity ($k$ Categorical Variables)}
Sparsity controls the number of categorical variables used to generate treatment-effect structure while keeping the total feature dimension fixed. Lower $k$ implies fewer variables with more categories per variable, which matches many third-party segmentation settings. DivergenceTree's recall advantage is largest at $k=1$ and $k=2$, while all methods converge once $k \ge 3$.
\begin{figure}[t!]
    \centering
    \begin{subfigure}[t]{0.49\textwidth}
        \centering
        \includegraphics[width=\textwidth]{figures/figurefiles/factor_comparison_sparsity_accuracy.png}
        \caption{Accuracy}
    \end{subfigure}
    \hfill
    \begin{subfigure}[t]{0.49\textwidth}
        \centering
        \includegraphics[width=\textwidth]{figures/figurefiles/factor_comparison_sparsity_recall_region_2.png}
        \caption{Recall for Region~2}
    \end{subfigure}
    \caption{\label{fig:factor-sparsity}
    Performance across sparsity levels ($k$ categorical variables).}
\end{figure}


\subsection{Rareness (Observation Share in Region~2)}
Rareness controls class imbalance for the diagnostic target. When Region~2 is rare, two-step recall drops sharply and can approach zero in the lowest bins, whereas DivergenceTree with $\lambda>0$ retains nontrivial detection power. This is where co-movement emphasis is most valuable.
\begin{figure}[t!]
    \centering
    \begin{subfigure}[t]{0.49\textwidth}
        \centering
        \includegraphics[width=\textwidth]{figures/figurefiles/factor_comparison_rareness_accuracy.png}
        \caption{Accuracy}
    \end{subfigure}
    \hfill
    \begin{subfigure}[t]{0.49\textwidth}
        \centering
        \includegraphics[width=\textwidth]{figures/figurefiles/factor_comparison_rareness_recall_region_2.png}
        \caption{Recall for Region~2}
    \end{subfigure}
    \caption{\label{fig:factor-rareness}
    Performance across rareness levels (share of observations in Region~2).}
\end{figure}


\subsection{Dispersion (Combination Share in Region~2)}
Dispersion controls how broadly contradiction combinations are distributed in feature space. Low dispersion means Region~2 is concentrated in a narrow set of combinations, making it harder to recover. DivergenceTree outperforms two-step most clearly in this low-dispersion regime.
\input{figures/fig_factor_dispersion}

\subsection{Noise (Treatment-Effect Noise Standard Deviation)}
Noise controls the signal-to-noise ratio of treatment effects. Co-movement emphasis delivers the largest recall gains in low-noise settings, where amplified joint deviations reflect true structure. As noise increases, all methods degrade and the incremental benefit of larger $\lambda$ diminishes.
\input{figures/fig_factor_noise}

\subsection{Data Size (Training Sample Size)}
Data size controls statistical power for segment recovery. At very small samples, two-step can be competitive due to bagging in causal forests. As sample size grows, DivergenceTree strengthens and maintains a clear recall advantage while preserving accuracy.
\begin{figure}[t!]
    \centering
    \begin{subfigure}[t]{0.49\textwidth}
        \centering
        \includegraphics[width=\textwidth]{figures/figurefiles/factor_comparison_data_size_accuracy.png}
        \caption{Accuracy}
    \end{subfigure}
    \hfill
    \begin{subfigure}[t]{0.49\textwidth}
        \centering
        \includegraphics[width=\textwidth]{figures/figurefiles/factor_comparison_data_size_recall_region_2.png}
        \caption{Recall for Region~2}
    \end{subfigure}
    \caption{\label{fig:factor-data-size}
    Performance across training sample sizes.}
\end{figure}



\end{appendices}
\end{bibunit}
\end{document}
